\chapter{Estado del Arte}\label{chapter:state-of-the-art}

Este capítulo revisa el estado actual de los sistemas de Generación Aumentada por Recuperación basados en grafos de conocimiento. El énfasis se pone en cómo estos sistemas gestionan la dimensión temporal. El objetivo es identificar qué problemas han sido resueltos, cuáles permanecen abiertos y qué decisiones de diseño justifican la propuesta del capítulo siguiente.

La revisión se organiza en siete secciones y un cierre. La sección~\ref{sec:graphrag-foundations} describe GraphRAG y sus limitaciones frente a corpus que evolucionan en el tiempo. La sección~\ref{sec:tkg-representations} examina las representaciones y los métodos de aprendizaje sobre grafos de conocimiento temporales. La sección~\ref{sec:temporal-reasoning} analiza el razonamiento temporal sobre grafos evolutivos a partir del sistema TREK. La sección~\ref{sec:temporal-graphrag} cubre los sistemas que incorporan componentes temporales al \emph{pipeline} de GraphRAG. La sección~\ref{sec:bitemporal} introduce el modelado bitemporal y los sistemas que lo han llevado a la práctica. La sección~\ref{sec:dynamic-update} aborda la actualización dinámica del grafo en dos vertientes: la resolución de entidades y el \emph{unlearning}. La sección~\ref{sec:benchmarks} discute los \emph{benchmarks} disponibles para evaluar este tipo de sistemas. El capítulo cierra en la sección~\ref{sec:soa-conclusions} con un resumen comparativo y con las limitaciones que motivan la propuesta de esta investigación.

%% ============================================================
\section{Generación Aumentada por Recuperación Basada en Grafos de Conocimiento}\label{sec:graphrag-foundations}

La Generación Aumentada por Recuperación (RAG) \cite{lewis2020rag} es el paradigma dominante para fundamentar las respuestas de LLMs en evidencia documental. En su formulación más simple, el sistema fragmenta el corpus en registros textuales y cada fragmento se codifica como un vector semántico. Ante cada consulta, el sistema recupera los fragmentos más relevantes y los incorpora como contexto en el \emph{prompt} del modelo. Este enfoque se denomina \emph{vector RAG} y resulta adecuado cuando la respuesta se localiza en unos pocos fragmentos. Pero es insuficiente para preguntas que exigen comprensión global del corpus, como ``¿cuáles son las tendencias principales del \emph{dataset}?'' \cite{edge2024local}.

GraphRAG (Graph Retrieval Augmented Generation, Microsoft Research) \cite{edge2024local} es un sistema de RAG que mejora las respuestas de los LLMs recuperando información de un corpus de documentos. Resuelve las limitaciones anteriores construyendo un grafo de conocimiento que captura la estructura relacional del corpus. Su \emph{pipeline} opera en dos fases:

\paragraph{Fase 1 -- Indexación.} En esta fase \cite{edge2024local}, los documentos se segmentan en bloques de texto. Un LLM extrae entidades y relaciones de cada bloque mediante \emph{prompts} con \emph{few-shot} (plantillas con unos pocos ejemplos resueltos que guían al modelo). La información que aparece varias veces en el corpus se consolida y las descripciones de las entidades se agregan en nodos únicos. Las relaciones duplicadas se fusionan como aristas ponderadas por frecuencia. La identificación de entidades equivalentes se realiza por coincidencia exacta del nombre extraído. Sobre el grafo resultante, el algoritmo de Leiden \cite{traag2019leiden} particiona los nodos en comunidades jerárquicas anidadas. Un LLM genera resúmenes descriptivos para cada comunidad en cada nivel, y así se construye una pirámide de descripciones de comunidades.
%TODO: figura de la pirámide de comunidades
GraphRAG extrae también afirmaciones factuales asociadas a entidades (\emph{claims}) \cite{edge2024local}, los únicos elementos con capacidad temporal en todo el sistema. El artículo original de GraphRAG describe los \emph{claims} como afirmaciones textuales factuales sobre entidades (fechas, eventos, interacciones), pero no formaliza un esquema estructurado. Es la implementación pública de Microsoft GraphRAG la que define esos \emph{claims} con campos tabulados. Los más relevantes para esta investigación son \texttt{start\_date} y \texttt{end\_date}, que delimitan el periodo de validez, junto con un \texttt{status} de veracidad.
%TODO: figura que ilustra la estructura de los claims

\paragraph{Fase 2 -- Consulta (\emph{query}).} El artículo original \cite{edge2024local} presenta la \emph{Búsqueda Global}, que se realiza mediante un esquema \emph{map-reduce}. El LLM resume cada comunidad por separado (\emph{map}) y después combina los resúmenes parciales en una respuesta única (\emph{reduce}). Esto permite la síntesis completa del corpus. La implementación de código abierto de GraphRAG incorpora también la \emph{Búsqueda Local}, que explora el vecindario de entidades relevantes, y la \emph{Búsqueda DRIFT}, que parte de una búsqueda global para situar el contexto y lo refina iterativamente con búsquedas locales sobre los nodos seleccionados. Estas dos modalidades no aparecen en el artículo original, sino solo en la implementación pública del repositorio de GraphRAG de Microsoft.

Los resultados experimentales del artículo \cite{edge2024local} se obtienen sobre dos corpus: \emph{podcasts} (${\sim}1$M \emph{tokens}, 8\,564 nodos, 20\,691 aristas) y noticias (${\sim}1.7$M \emph{tokens}, 15\,754 nodos, 19\,520 aristas). GraphRAG supera al RAG vectorial en capacidad de comprensión (72--83\% en \emph{podcasts} y 72--80\% en noticias). También lo supera en diversidad de las respuestas (75--82\% y 62--71\%, respectivamente). Ambas métricas se evalúan con un juez LLM que compara pares de respuestas y emite una preferencia. Cuando la consulta se responde desde el nivel más alto de la jerarquía de comunidades (es decir, usando solo los resúmenes más generales del corpus), GraphRAG necesita más de un 97\% menos de \emph{tokens} de contexto que el \emph{baseline} ingenuo. Ese \emph{baseline} consiste en pasar el corpus completo al LLM en un único \emph{prompt} y pedirle un resumen (sumarización directa).

Cuando el corpus subyacente evoluciona en el tiempo, \textbf{GraphRAG carece de mecanismos para gestionar esa evolución a nivel estructural del grafo}. Ni las entidades, ni las relaciones, ni las descripciones de las comunidades contienen información sobre su periodo de vigencia. Por ejemplo, una relación ``dirige'' registrada a partir de un documento de 2010 resulta indistinguible de otra extraída de un documento de 2025.

Otro aspecto a resaltar es que los \emph{claims} no se usan para estructurar el grafo temporalmente. Solo se insertan como anotaciones textuales en el \emph{prompt} del LLM, no actúan como restricciones para filtrar la recuperación y tampoco condicionan la sumarización del corpus.

Además de los aspectos antes señalados, se observa que los \emph{claims} capturan solo el \emph{tiempo válido} (cuándo el hecho fue verdadero en el mundo real). Omiten el \emph{tiempo de transacción} (cuándo el sistema incorporó ese hecho). Esto impide diferenciar dos casos: una \textbf{corrección retroactiva} es un hecho que en un momento dado se asumió como correcto y luego se evidenció que no lo era; una \textbf{evolución natural} es un hecho que fue verdadero durante un periodo y luego dejó de serlo. Ambas nociones se formalizan en la sección~\ref{sec:bitemporal}.

A estas limitaciones se añaden dos problemas operativos que el artículo original no discute. El primero: el sistema permite la coexistencia de hechos contradictorios sin detectarlos. Por ejemplo, si un documento de 2019 y otro de 2023 atribuyen el mismo cargo directivo a dos personas distintas, ambos permanecen en el grafo sin alerta de conflicto.

El segundo: la coincidencia exacta descrita en el artículo de GraphRAG \cite{edge2024local} fragmenta la resolución de entidades ante variaciones nominales. ``Organización de las Naciones Unidas'' y ``ONU'' generan nodos separados. Las relaciones que deberían referirse a la misma entidad ni siquiera se conectan en el grafo.

Estas insuficiencias estructuran los siguientes acápites del capítulo: la ausencia de temporalidad como dimensión estructural, la pasividad de los \emph{claims}, la acumulación de contradicciones sin resolución y la fragilidad en la resolución de entidades. El análisis se centrará en ellas para justificar la idoneidad de la propuesta de solución.

%% ============================================================
\section{Grafos de Conocimiento Temporales: Representaciones y Aprendizaje}\label{sec:tkg-representations}

Los Grafos de Conocimiento Temporales (\emph{Temporal Knowledge Graphs}, TKGs) \cite{tkg_survey2024} extienden los grafos de conocimiento estáticos. Incorporan la dimensión temporal en la representación de los hechos. En un grafo estático, un hecho se considera verdadero de manera atemporal. Afirmar que una persona dirige una empresa no indica desde cuándo ni hasta cuándo. Los TKGs parten de una premisa distinta: la mayoría del conocimiento del mundo real solo es válido en periodos específicos y evoluciona con el tiempo. Por tanto, la representación debe incorporar la temporalidad como elemento de primera clase.

Formalmente, un TKG se define como $G = (E, R, T, F)$ \cite{tkg_survey2024}. $E$ es el conjunto de entidades, $R$ es el de relaciones, $T$ es el de marcas temporales y $F \subset E \times R \times E \times T$ es el conjunto de hechos. Esta definición extiende la representación clásica de tripletas $(s, r, o)$, donde $s$ es el sujeto, $r$ la relación y $o$ el objeto. La extensión añade una componente temporal. En la práctica, esa componente adopta dos formas principales.

La primera son las \emph{cuádruplas} $(s, r, o, t)$ \cite{tkg_survey2024}. Asocian cada hecho a un punto temporal discreto ($t$). Por ejemplo, (\textit{China}, \textit{organizó}, \textit{Juegos Olímpicos de Verano}, 2008-08) indica que China organizó los Juegos en agosto de 2008. No aporta información sobre su duración. La segunda son las \emph{quíntuplas} $(s, r, o, t_{\text{start}}, t_{\text{end}})$, donde $t_{\text{start}}$ y $t_{\text{end}}$ son el inicio y el fin del intervalo de validez. Se corresponden con los \emph{rangos temporales} (\emph{range timestamps}) descritos por \cite{tkgc_survey2023}. Por ejemplo, (\textit{Kofi Annan}, \textit{Secretario\_General\_de}, \textit{ONU}, 1997-01, 2006-12) indica que Kofi Annan ocupó el cargo durante un periodo concreto.

%TODO: figura fig:tkg-cuadruplas-quintuplas que compare visualmente la representación de un hecho como cuádrupla (s, r, o, t) con marca temporal puntual frente a su representación como quíntupla (s, r, o, t_start, t_end) con intervalo de validez, usando el ejemplo de Kofi Annan como Secretario General de la ONU.
La distinción entre ambas representaciones tiene consecuencias prácticas. Las cuádruplas son adecuadas para eventos puntuales. Las quíntuplas capturan estados con duración. Las quíntuplas son empleadas por bases de conocimiento como Wikidata o YAGO \cite{tkgc_survey2023}. Desde el punto de vista del razonamiento, las representaciones puntuales limitan las consultas a coincidencias exactas de \emph{timestamp} (marca que identifica un instante de tiempo). Las que tiene el intervalo de tiempo válido habilitan consultas sobre solapamiento, contención y precedencia temporal \cite{tkg_survey2024, tkgc_survey2023}. Esto amplía la expresividad del lenguaje de consulta.

Sobre estas representaciones, la literatura refleja un grupo amplio de métodos de aprendizaje. Los \emph{surveys} recientes \cite{tkg_survey2024, tkgc_survey2023} los sistematizan en familias diferenciadas. La diferencia principal está en la estrategia para incorporar el tiempo.

\paragraph{Métodos basados en transformaciones.} Extienden modelos como TransE \cite{bordes2013translating}. Codifican las marcas temporales como traslaciones o rotaciones adicionales en el espacio de \emph{embeddings}. Son representativos TTransE, HyTE y ChronoR \cite{tkg_survey2024}.

\paragraph{Métodos de descomposición tensorial.} Tratan al TKG como un tensor de cuatro dimensiones (sujeto $\times$ relación $\times$ objeto $\times$ tiempo). Estos métodos aprenden vectores de baja dimensión a partir del conjunto de entrenamiento. El producto de esos vectores reconstruye el tensor original. TComplEx \cite{tkg_survey2024} y TNTComplEx \cite{tkgc_survey2023} son representativos.

\paragraph{Métodos basados en GNN y recurrencia.} Combinan redes neuronales sobre grafos (GNN, \emph{Graph Neural Networks} \cite{scarselli2009gnn}) con módulos recurrentes. Son representativos RE-NET \cite{jin2020renet} y RE-GCN \cite{li2021regcn}. Estos procesan el TKG como una secuencia de \emph{snapshots} discretos, donde cada \emph{snapshot} es una fotografía del grafo con los hechos válidos en un intervalo fijo. La información temporal se propaga mediante unidades recurrentes sobre una base R-GCN (\emph{Relational Graph Convolutional Network} \cite{schlichtkrull2018rgcn}). Las unidades recurrentes empleadas son GRU (\emph{Gated Recurrent Unit} \cite{chung2014gru}) y LSTM (\emph{Long Short-Term Memory} \cite{hochreiter1997lstm}).

\paragraph{Métodos basados en procesos puntuales.} Modelan los hechos como eventos discretos en tiempo continuo. Cada evento se define por una función de intensidad. La función predice la probabilidad instantánea de que ocurra un nuevo evento. Son representativos Know-Evolve, GHNN y EvoKG \cite{tkg_survey2024}.

\paragraph{Métodos basados en hipergrafos.} DHE-TKG \cite{dhe_tkg2024} extiende las representaciones a estructuras de orden superior. Emplea hipergrafos. Esos hipergrafos agrupan entidades relacionadas en cada \emph{timestamp}.

Los \emph{surveys} \cite{tkg_survey2024, tkgc_survey2023} reportan que los autores de cada método evalúan sus propuestas sobre dos tipos de \emph{datasets}. Por un lado, \emph{datasets} de eventos políticos extraídos de noticias: ICEWS~\cite{boschee2015icews} y GDELT~\cite{leetaru2013gdelt}. Por otro, bases de conocimiento de propósito general: Wikidata~\cite{vrandecic2014wikidata} y YAGO~\cite{suchanek2007yago}. Las métricas estándar son dos. \emph{Mean Reciprocal Rank} (MRR) \cite{manning2008ir} es el promedio del inverso del rango de la respuesta correcta. Hits@$k$ \cite{bordes2013translating} es la fracción de consultas en las que la respuesta correcta está entre las $k$ mejores \cite{tkg_survey2024, tkgc_survey2023}.

Como balance del panorama anterior, la comunidad ha desarrollado mecanismos formalmente maduros para modelar la evolución del conocimiento. Hay funciones de decaimiento temporal, que reducen el peso de los hechos antiguos. Hay dependencias secuenciales mediante módulos autorregresivos, que predicen el siguiente paso a partir de los anteriores. Y hay tratamiento continuo del tiempo mediante procesos puntuales \cite{tkg_survey2024, tkgc_survey2023}.

Sin embargo, todos los métodos revisados operan sobre un único eje temporal. Ese eje se denomina \emph{tiempo válido} (\emph{valid time}) \cite{tkg_survey2024, tkgc_survey2023} y registra cuándo un hecho es verdadero en el mundo real. Los \emph{surveys} \cite{tkg_survey2024, tkgc_survey2023} señalan tres direcciones futuras: la escalabilidad, la interpretabilidad y la integración con modelos de lenguaje de gran escala. El \emph{survey} de \emph{Temporal Knowledge Graph Completion} \cite{tkgc_survey2023} menciona además el \emph{catastrophic forgetting} en escenarios de aprendizaje continuo. El \emph{catastrophic forgetting} es la tendencia de los modelos neuronales a perder lo aprendido sobre datos antiguos cuando se les entrena con datos nuevos. Sin embargo, ninguno de los dos \emph{surveys} identifica como problema abierto la capacidad del modelo de eliminar el efecto de un dato concreto, de modo que el resultado sea indistinguible del que se obtendría volviendo a ejecutar el sistema desde cero sin ese dato (\textbf{\emph{unlearning}}). Tampoco identifican la capacidad de detectar y resolver contradicciones entre hechos de diferentes periodos (\textbf{resolución de conflictos temporales}).

Los problemas anteriores tienen un factor común. Al registrar únicamente el tiempo válido, los TKGs no distinguen entre cuándo un hecho es verdadero y cuándo el sistema tomó conocimiento de ese hecho. Esa omisión impide diferenciar dos escenarios distintos. El primero es una \textbf{corrección retroactiva}. Por ejemplo, en 2022 el sistema descubre que un hecho que creía verdadero desde 2020 era incorrecto (un dato mal extraído por el LLM). El segundo es una \textbf{evolución natural}. Por ejemplo, un hecho que era verdadero en 2020 dejó de serlo en 2023 (un cambio de CEO en una compañía). Ambas situaciones requieren operaciones de actualización diferentes. El único eje temporal disponible las colapsa en la misma representación.

A esto se añade otro problema. El enfoque predominante se basa en \emph{snapshots} discretos \cite{tkg_survey2024}. Es adecuado para los \emph{datasets} estáticos empleados en la evaluación, pero resulta insuficiente cuando el conocimiento llega de manera incremental y continua, que es el escenario de interés para un sistema de GraphRAG temporal dinámico.

%% ============================================================
\section{Razonamiento Temporal sobre Grafos Evolutivos}\label{sec:temporal-reasoning}

TREK (\emph{Temporal Reasoning on Evolving Knowledge Graphs}) \cite{trek2025} aborda los dos aspectos que quedaron abiertos en la sección anterior. Estos aspectos son la gestión de conflictos y la distinción entre corrección y evolución.

La propuesta tiene dos componentes. EVOKG \cite{trek2025} se ocupa de cómo evoluciona el grafo cuando llega información nueva. EVOREASONER \cite{trek2025} se ocupa de cómo se razona sobre el grafo ante consultas temporalmente complejas. La separación responde a que actualizar y consultar conocimiento son problemas que requieren un enfoque distinto.

El núcleo de EVOKG es una clasificación de las relaciones según su cardinalidad temporal. Una relación es \emph{exclusiva} si una entidad solo admite una instancia activa en un momento dado. Ejemplos son la fecha de nacimiento de una persona o la capital de un país. Una relación es \emph{no exclusiva} si admite varias instancias activas a la vez. Ejemplos son los idiomas que habla una persona o los libros que ha publicado un autor. Esta distinción condiciona la estrategia de actualización. Para las relaciones no exclusivas, el sistema añade el nuevo hecho al historial con su propio intervalo de validez. En las exclusivas, la llegada de un hecho nuevo plantea un conflicto. TREK resuelve este conflicto con una fórmula de confianza que pondera tres señales. La primera es la frecuencia con que el hecho aparece en las fuentes. La segunda es cuán reciente es el hecho, que se modela con una ponderación sigmoidal $1/(1+e^{-\gamma \Delta t})$ sobre el tiempo transcurrido $\Delta t$~\cite{trek2025}. La tercera es el \emph{peso} atribuido a cada fuente.

EVOREASONER procesa las consultas que no se resuelven con una sola búsqueda. Un ejemplo es ``¿quién era el director de X cuando se fusionó con Y?'', que encadena dos sub-consultas. Primero hay que determinar la fecha de la fusión. Después hay que recuperar el director vigente en esa fecha. EVOREASONER descompone estas preguntas en rutas de exploración del grafo. Cada ruta es una cadena de relaciones sujeta a restricciones temporales. El sistema usa una búsqueda tipo \emph{beam}~\cite{lowerre1976harpy}, que mantiene en cada paso los $n$ caminos parciales con mayor puntuación. El \emph{scoring} es sensible al tiempo. Penaliza los caminos cuyas relaciones tienen marcas temporales incompatibles con la consulta o demasiado alejadas en el tiempo.

Los resultados experimentales sostienen la utilidad del enfoque. La evaluación se realiza en dos escenarios. El primero es KGQA (\emph{Knowledge Graph Question Answering})~\cite{lan2022complex} puro sobre los \emph{benchmarks} TimeQuestions~\cite{jia2021timeq} y MultiTQ~\cite{chen2023multitq}. En estos \emph{benchmarks} TREK alcanza mejoras absolutas de hasta el 23.3\% y el 18.1\% en \emph{accuracy} respectivamente. Los \emph{baselines} cubren \emph{prompting} directo y razonamiento paso a paso (\emph{Chain-of-Thought}~\cite{wei2022cot}, \emph{Self-Consistency}~\cite{wang2023selfconsistency}). Cubren también variantes basadas en recuperación: RAG y RAG ampliado con los vecinos directos en el KG de las entidades recuperadas. Y abarcan métodos multi-salto sobre grafos como Think-on-Graph~\cite{sun2024tog} y Plan-on-Graph~\cite{luo2024pog}. El segundo escenario es \emph{end-to-end} con evolución del KG sobre CRAG~\cite{crag2024}, en sus dominios de películas y deportes. Aquí EVOKG procesa unos 3\,400 episodios de actualización en 24 horas. Un modelo de 8 mil millones de parámetros, del orden de magnitud de Llama-3-8B~\cite{llama3_2024}, equipado con TREK iguala a un modelo de 671 mil millones, DeepSeek-V3~\cite{deepseekv3_2025}, entrenado siete meses después y consultado por \emph{prompting} puro sin EVOKG. La gestión activa del tiempo compensa dos órdenes de magnitud de capacidad paramétrica y, en parte, la ventaja de un entrenamiento más reciente.

%TODO: figura fig:trek-cardinalidad-asimetrica que ilustre la asimetría exclusiva/no-exclusiva con dos ejemplos (``es\_CEO\_de'' y ``nació\_en'') mostrando, mediante diagramas de aristas, cómo la restricción opera desde un extremo de la relación pero no desde el otro.
En este trabajo se ha identificado una limitación de la clasificación de relaciones de TREK que el propio artículo no discute. La etiqueta exclusiva/no-exclusiva se asigna a la relación como un todo. No distingue desde qué extremo de la relación opera la restricción. Considérese la relación ``es\_CEO\_de''. Vista desde la empresa es exclusiva, porque una empresa tiene un único CEO a la vez. Vista desde la persona deja de serlo, porque una misma persona puede dirigir varias empresas en paralelo. La relación ``nació\_en'' presenta la asimetría inversa. Vista desde la persona es exclusiva, porque una persona nace en un único lugar. Vista desde el lugar deja de serlo, porque en una misma ciudad pueden haber nacido múltiples personas. Una etiqueta única no captura este matiz. La consecuencia práctica es directa. El camino que se sigue en la búsqueda de un conflicto depende del extremo que impone la restricción. Clasificar de forma simétrica genera falsos positivos y deja pasar conflictos reales.

TREK formaliza la distinción entre reemplazo y corrección que los TKGs anteriores no abordaban. Demuestra que la evolución del conocimiento puede resolverse de manera dinámica y no solo acumulativa. Su clasificación de las relaciones en exclusivas y no exclusivas resulta, sin embargo, insuficiente, y plantea la cuestión de cómo modelar relaciones con restricción asimétrica respecto a sus extremos. TREK opera además sobre TKGs generales y no sobre el \emph{pipeline} de GraphRAG. La pregunta inmediata es cómo trasladan estas intuiciones los sistemas construidos sobre dicho \emph{pipeline}. La solución que este trabajo propone para la asimetría de las restricciones de cardinalidad se desarrolla en la sección ~\ref{chapter:proposal}.

%% ============================================================
\section{Sistemas de GraphRAG con Conciencia Temporal}\label{sec:temporal-graphrag}

Esta sección examina los sistemas que han incorporado componentes temporales al \emph{pipeline} de GraphRAG o a arquitecturas afines de RAG basado en grafos. El orden que se estructura esta sección se corresponde con la anatomía del \emph{pipeline} es decir: representación del tiempo en el grafo, descomposición y resolución de consultas y estrategias ante hechos contradictorios.

\subsection{Representación del tiempo en el grafo}

El modelo DyG-RAG \cite{dyg_rag2025} introduce un cambio conceptual con respecto al tratamiento simple de la temporalidad que propone GraphRAG. El mismo consiste en hacer del evento, y no de la entidad, la unidad atómica de representación. Cada nodo del grafo es una \emph{Dynamic Event Unit} (DEU)\cite{dyg_rag2025}el cual lleva asociado un \emph{timestamp} explícito. Las aristas del grafo se ponderan combinando dos factores: la co-ocurrencia de las entidades mencionadas en cada evento y la proximidad temporal entre ellos. Este diseño propuesto centrado en los DEU tiene dos consecuencias relevantes: el tiempo deja de ser un atributo pasivo y juega un papel importante en la puntuación de relevancia durante la recuperación. De igual forma la inserción incremental de nuevos eventos resulta natural a partir del diseño propuesto. La representación se complementa con una estrategia de \emph{prompting} denominada \emph{Time Chain-of-Thought} (Time-CoT)\cite{dyg_rag2025}, que obliga al LLM a reconstruir explícitamente la línea temporal de los eventos antes de generar la respuesta. Sobre los \emph{benchmarks} TimeQA~\cite{chen2021timeqa}, TempReason~\cite{tan2023tempreason} y ComplexTR~\cite{tan2023complextr}, todos derivados de Wikipedia, DyG-RAG reporta ganancias de \emph{accuracy} de entre un 10\% y un 18\% frente a GraphRAG~\cite{edge2024local}, LightRAG~\cite{guo2024lightrag} e HippoRAG~\cite{gutierrez2024hipporag}, empleando 4 GPUs NVIDIA V100.

Por otra parte, el modelo E2RAG aborda un problema que otros sistemas pasan por alto: fusionar menciones de una misma entidad sin tener en cuenta el tiempo elimina información relevante.  La solucion que brinda para resolver dicho problema parte desde el propio modelado de los datos. Consiste en construir dos grafos paralelos donde los nodos de los mismos se conectan mediante aristas bipartitas. En uno grafo se almacena las menciones de entidades y el otro almacena los eventos. Cada mención del texto origina un nodo nuevo, conectado al evento donde aparece. Así, dos menciones de una misma entidad producen dos nodos distintos y la resolución de entidades queda aplazada hasta el momento de la consulta.

Por ejemplo, en uno de los libros de Harry Potter el personaje se menciona en el capítulo 1 y en el capítulo 7. E2RAG crea un nodo distinto para cada mención del personaje, mientras que los sistemas tradicionales fusionan el personaje en una misma entidad y por tanto pierden su evolución entre ambos capítulos.

En la evaluación sobre el \emph{benchmark} ChronoQA \cite{e2rag2025}, compuesto por 497 pares pregunta-respuesta extraídos de 9 novelas, E2RAG obtiene la mejor puntuación entre los sistemas comparados en \cite{e2rag2025} .

Las dos propuestas que se están analizando convierten el tiempo en un factor activo del modelado: DyG-RAG lo integra en el \emph{scoring} y E2RAG lo usa para evitar fusionar menciones de una misma entidad en momentos distintos. Sin embargo, las dos representan el tiempo como puntos discretos y no como intervalos de validez, por lo que no distinguen entre hechos caducados de hechos falsos. Además DyG-RAG tampoco detecta conflictos entre eventos contradictorios, y la estrategia de E2RAG de no fusionar nunca las menciones es eficaz en narrativas autocontenidas pero se degrada en corpus extensos por la fragmentación de entidades.

\subsection{Descomposición y resolución de consultas}

Otros investigaciones analizadas desplazan el énfasis desde la representación hacia la consulta. Por ejemplo en MRAG \cite{mrag2025} se descompone cada pregunta en dos componentes: el contenido semántico y la restricción temporal. En el \emph{ranking} usado por este se combina similitud vectorial para el contenido con una puntuación temporal para la restricción. La separación reporta mejoras de ${\sim}9.3\%$ en \emph{answer recall} (fracción de preguntas cuya respuesta correcta aparece entre los fragmentos recuperados) y de ${\sim}4.5\%$ en EM/F1 (\emph{Exact Match}, coincidencia textual exacta con la respuesta de referencia, y \emph{F1}, solapamiento a nivel de \emph{tokens}).

En \cite{kg_irag2025} donde se aborda el modelo KG-IRAG se le da prioridad a la representación del conocimiento en lugar de la consulta. El modelo representa las marcas temporales (\emph{timestamps}) como nodos del propio grafo, conectando a todos los hechos con sus marcas temporales. La consulta se resuelve mediante un proceso iterativo Planner-Judge~\cite{kg_irag2025}. El Planner se encarga de identificar un intervalo y una entidad iniciales y genera el \emph{prompt} de razonamiento. Por su parte el Judge evalúa si la información recuperada basta para responder; si no, decide qué nuevo intervalo, o entidad, consultar.

Ambos sistemas confirman que separar la dimensión temporal mejora la calidad de las respuestas. Sin embargo, comparten una limitación de fondo: operan sobre grafos estáticos, construidos una sola vez. KG-IRAG presenta además una decisión de diseño cuestionable: al representar los periodos como nodos y no como atributos de las aristas, confunde las entidades del dominio con simples marcas temporales. Esta confusión entorpece las consultas que cruzan periodos de tiempo y mezcla la dimensión semántica con la temporal.

\subsection{Estrategias ante hechos contradictorios}

Otros sistemas abordan de forma explícita la coexistencia de hechos potencialmente contradictorios. En T-GRAG \cite{tgrag2025} se parte de un módulo del sistema denominado \emph{Temporal Query Decomposer} que divide cada consulta multi-temporal en sub-consultas, cada una asociada a una única restricción temporal. La recuperación de la información representada en el grafo se organiza en tres capas: un \emph{Temporal Subgraph Retriever} que restringe el grafo al periodo relevante; un \emph{Coarse Node Retrieval} que selecciona nodos por similitud semántica dentro de ese subgrafo seleccionado; y un \emph{Fine Knowledge Retrieval} que extrae del mismo los hechos concretos junto con sus \emph{timestamps}.

% TODO: figura fig:tgrag-three-layer-retrieval que muestre las tres capas de recuperación de T-GRAG (Temporal Subgraph Retriever, Coarse Node Retrieval y Fine Knowledge Retrieval) como un diagrama vertical con el flujo de datos entre ellas y la consulta de entrada.

En la evaluación del modelo anterior sobre el dataset Time-LongQA (2\,292 pares pregunta-respuesta extraídos de los informes anuales de Audi entre 2012 y 2023), T-GRAG supera a GraphRAG entre 34 y 47 puntos porcentuales de \emph{accuracy} en preguntas con varias referencias temporales.

La mejora es notable, pero T-GRAG no resuelve los conflictos; los \emph{evita} a partir del diseño asumido. Guarda cada hecho con su \emph{timestamp} y atiende cada sub-consulta dentro de su periodo de tiempo correspondiente. Así, delega toda la coherencia de la respuesta en el filtrado temporal de la consulta, sin detectar ni resolver las contradicciones temporales presentes en el grafo.

%TODO: figura fig:tgrag-bi-nivel que ilustre la arquitectura bi-nivel de TG-RAG: en el nivel inferior, el grafo de entidades con aristas etiquetadas con timestamps; en el nivel superior, la jerarquía de nodos temporales (año, trimestre, mes, día) con un time report asociado a cada uno; conexiones verticales entre niveles indicando la cobertura temporal.
Por su parte en TG-RAG \cite{tg_rag2025} se avanza en la misma dirección con una arquitectura de grafo temporal bi-nivel. El nivel inferior contiene los nodos referentes a entidades enlazadas por aristas con un \emph{timestamps} asociado. El nivel superior organiza nodos temporales en una jerarquía creciente entre dichos nodos (año, trimestre, mes, día). Para cada nodo temporal, el LLM genera un \emph{time report} que sintetiza el estado del conocimiento en el periodo de tiempo correspondiente.

TG-RAG alcanza una \emph{factuality} (proporción de afirmaciones generadas que son consistentes con la información del corpus) de 0.599 frente a aproximadamente 0.410 de los \emph{baselines} sobre el corpus ECT-QA \cite{tg_rag2025}. La aportación más significativa del sistema es la viabilidad de la actualización incremental del corpus. Cuando llegan nuevos documentos basta con regenerar los \emph{time reports} de los nodos temporales nuevos y de sus ancestros. 

En el tratamiento de los conflictos, TG-RAG sigue el mismo patrón que T-GRAG. Los hechos de momentos distintos se almacenan como aristas (arreglarlo), sin detección explícita de contradicciones. El filtrado temporal en la consulta separa los hechos por periodo de tiempo.

La Tabla~\ref{tab:state-of-the-art-comparison} contiene el análisis comparativo de los modelos T-GRAG y TG-RAG en cuanto su diseño. A partir de dicha compración se determina que ambos sistemas adoptan una estrategia que evade la resolución de conflictos temporales.


Los sistemas revisados muestran un progreso en la incorporación de componentes temporales al \emph{pipeline} de GraphRAG. El análisis resumido en la Tabla~\ref{tab:state-of-the-art-comparison} revela que la temporalidad se aborda en la fase de consulta, mediante filtrado, descomposición y \emph{scoring} temporales pero no en la fase de indexación. Al no tener en cuenta la temporalidad en la fase de indexación, los sistemas revisados no detectan ni resuelven los conflictos temporales presentes en el grafo a medida que el corpus evoluciona dinamicamente a traves del teimpo.
%% ============================================================
\section{Modelado Bitemporal: Fundamentos Teóricos y Sistemas Existentes}\label{sec:bitemporal}

Los sistemas examinados hasta aquí modelan la temporalidad solo a partir del tiempo en que un hecho es verdadero en el mundo real y omiten el momento en que el sistema llegó a conocer dicho hecho. Esa omisión vuelve indistinguible la corrección retroactiva de la evolución natural explicadas en la sección~\ref{sec:graphrag-foundations} del presente documento. Esto hace que un documento tardío sobre un hecho antiguo aparezca en el grafo como si siempre hubiera estado allí, y no permite diferenciar la información reciente de la antigua ya existente en el sistema. Esta sección revisa los fundamentos teóricos del modelado bitemporal y los sistemas que lo han llevado a la práctica, para evaluar su posible integración en un \emph{pipeline} de GraphRAG.

El marco bitemporal separa dos intervalos temporales independientes. El primero, denominado tiempo válido (\emph{valid time}), registra cuándo un hecho es verdadero en el mundo real. El segundo, denominado tiempo de transacción (\emph{transaction time}), registra cuándo el sistema tomó conocimiento de él.
%TODO: figura fig:bitemporal-cuadrantes que represente, en un plano cartesiano cuyo eje X sea el tiempo válido y el eje Y el tiempo de transacción, los cuatro estados cualitativos del modelado bitemporal (verdad actual, verdad histórica, error retractado y corrección retroactiva), con un ejemplo en cada cuadrante.
A partir de estos dos intervalos la información que brindan los mismos puede conllevar a cuatro estados cualitativamente distintos~---verdad actual, verdad histórica, error retractado y corrección retroactiva---~los cuales no pueden representarse si solo se contase con un único intervalo, como ocurre en los modelos explicados en las secciones anteriores.

En el plano formal, los lenguajes para manejar datos temporales TLINQ y VLINQ \cite{tlinq2022} formalizan por separado las consultas sobre el tiempo de transacción y sobre el tiempo válido. En los resultados expuestos en dicha publicación se destacan dos aportes: la técnica \emph{sequenced join}, a partir de la cual se automatiza la intersección de intervalos de validez y elimina una fuente histórica de errores, y la distinción formal entre actualizaciones secuenciadas (evolución) y no secuenciadas (corrección), con traducciones acompañadas de pruebas de corrección. Que ambos cálculos se desarrollen por separado deja planteada de forma implícita la unificación en un único cálculo bitemporal.

%TODO: figura fig:zep-graphiti-architecture que muestre los tres niveles de ZEP/Graphiti (Episode Subgraph, Semantic Entity Subgraph y Community Subgraph) como capas apiladas, con un detalle aparte de las cuatro marcas temporales asociadas a cada arista del nivel intermedio ($t_{\text{valid}}$, $t_{\text{invalid}}$, $t'_{\text{created}}$, $t'_{\text{expired}}$).
El sistema más cercano a un modelo bitemporal es ZEP/Graphiti \cite{zep2025}. Su arquitectura se organiza en tres niveles: un \emph{Episode Subgraph} que contiene los datos ``crudos''; un \emph{Semantic Entity Subgraph} que contiene las entidades y relaciones; y un \emph{Community Subgraph} basado en \emph{Dynamic Label Propagation} \cite{zhu2002labelpropagation}. Lo relevante para este trabajo es su modelo de datos a nivel de arista: donde cada una lleva cuatro marcas temporales, $t_{\text{valid}}$ y $t_{\text{invalid}}$ para el tiempo válido y $t'_{\text{created}}$ y $t'_{\text{expired}}$ para el tiempo de transacción. La detección de contradicciones temporales (expresadas por algunas aristas del grafo) se realiza mediante el uso de un LLM. En este proceso se identifican aristas a invalidar debido a las contradicciones detectadas, para luego invalidar cada arista fijando su parámetro $t_{\text{invalid}}$ a la fecha actual del procesamiento.

En \cite{zep2025} se presentan tres limitaciones relevantes que tiene el sistema ZEP/Graphiti. La primera de ellas es que está concebido como memoria conversacional para agentes de IA y no como instrumento para indexar corpus documentales. La segunda es que la detección de contradicciones se apoya en la similitud semántica entre aristas, sin una clasificación estructurada del tipo de relación ni de sus restricciones de cardinalidad, lo que le impide discernir si un conflicto procede de una relación exclusiva o de una que admite coexistencia legítima de instancias. La última de ellas es que no opera dentro del \emph{pipeline} de GraphRAG. Por tanto, la bitemporalidad y el \emph{pipeline} de GraphRAG permanecen en sistemas separados.

Como conclusión del análisis de los temas abordados en esta sección se puede plantear que ningún sistema ha integrado aún la semántica bitemporal dentro del \emph{pipeline} de GraphRAG. Este aspecto constituye una de las principales motivaciones para la realización de la presente investigación. Integrar la bitemporalidad a un grafo que evoluciona dinámicamente en el tiempo exige, como condición previa, resolver cómo se realizará la actualización del grafo de conocimiento ante la incorporación de nuevos documentos al corpus. Dicha actualización lleva implícita la solución de los siguientes aspectos:
\begin{itemize}
    \item cuándo una nueva mención en alguno de los documentos nuevos que se incorporan al corpus se corresponde con alguna entidad que ya está presente en la base de conocimientos;
    \item cuándo un hecho representado en la base de conocimientos debe retirarse de la misma por algún motivo relacionado con la actualización dinámica del corpus.
\end{itemize}
En la siguiente sección se analiza el estado actual de la técnica con respecto a estos dos problemas.

%% ============================================================
\section{Actualización Dinámica de Grafos: Resolución de Entidades, Unlearning y Corrección de Conocimiento}\label{sec:dynamic-update}

Esta sección aborda cómo actualizar un grafo de conocimiento cuando llega información nueva. El problema tiene tres aspectos importantes:
\begin{itemize}
    \item La resolución de entidades decide si una mención en un documento nuevo corresponde a una entidad ya existente o a una distinta.
    \item El \emph{unlearning} trata de cómo eliminar de una base de conocimiento algunos datos que han dejado de ser válidos.
\end{itemize}
A continuación, se abordan, en detalles, ambos aspectos.

\subsection{Resolución de entidades y construcción incremental de grafos}\label{subsec:entity-resolution}

En \cite{dalk2024} y \cite{dynamo2025} se aborda el problema de la resolución de entidades en sistemas que construyen grafos de forma incremental. En ambos reportes se presentan, y se analizan, dos sistemas de este tipo asistidos por LLMs: DALK \cite{dalk2024} y DYNAMO \cite{dynamo2025}.

DALK trabaja en el dominio biomédico, con preguntas sobre la enfermedad de Alzheimer extraídas de PubMed \cite{pubmed}. Propone una co-augmentación entre el LLM y el grafo. El grafo aporta contexto al modelo y el modelo extrae nuevas relaciones para añadir al grafo. La resolución de entidades se hace por \emph{entity linking} \cite{sevgili2022neural} (asociar a cada entidad una mención textual que describa el contexto del cual se extrajo la misma) basado en similitud semántica y haciendo uso de un modelo de \emph{sentence similarity} \cite{reimers2019sbert} (el modelo mide la cercanía de significado entre dos fragmentos de texto a partir de la distancia entre sus \emph{embeddings}). Como modelo principal de razonamiento usan GPT-3.5-turbo \cite{openai2023gpt4}, con evaluaciones adicionales sobre GPT-4 \cite{openai2023gpt4} y LLaMA \cite{llama3_2024}. DALK reporta mejoras absolutas de 16--19 puntos porcentuales sobre modelos biomédicos como Biomistral-7B \cite{labrak2024biomistral} y Meditron-70B \cite{chen2023meditron}.

DYNAMO se aplica a noticias verificadas y se evalúa sobre los \emph{datasets}: Hover \cite{jiang2020hover} (18\,171 muestras) y Feverous \cite{aly2021feverous} (26\,928 muestras). Este sistema introduce un principio que DALK no contempla: la verificación como criterio para incorporar información al grafo. Para la incorporación de información nueva al grafo y por ende, para la resolución de entidades, usa un módulo de \emph{fact-checking} \cite{guo2022survey} (verificación de hechos) basado en un LLM con razonamiento Chain-of-Thought \cite{wei2022cot} se combina con la detección de noticias falsas. Solo el conocimiento extraído de noticias verificadas como reales se incorpora al grafo. El principio básico en el cual se basa esta investigación es el siguiente: no toda información nueva debe incorporarse al grafo de conocimiento.

%TODO: figura fig:apple-temporal-resolution que ilustre el ejemplo ``Apple 1985 vs.\ 2024'': dos menciones de la misma entidad nominal en líneas de tiempo separadas por décadas, con descripciones contrastadas (empresa en crisis vs.\ compañía más valiosa) y la indicación visual de que un sistema sin temporalidad fusionaría ambos nodos perdiendo la diferencia de estado.
Tanto DALK como DYNAMO presentan una limitación: la resolución de entidades es puramente semántica y no tiene en cuenta la temporalidad. O sea, a partir de sus enfoques, no se contempla la posibilidad de que dos menciones de una misma entidad, en periodos distintos, puedan hacer referencia a estados muy diferentes de dicha entidad. Por ejemplo, ``Apple'' en 1985 era una empresa en crisis. ``Apple'' en 2024 es la compañía más valiosa del mundo. Es la misma entidad, pero al fusionar las menciones sin su contexto temporal hace que se pierda información relevante. En este caso el estado tan diferente en que se encuentra la empresa en 2024 con respecto a cómo se encontraba en 1985 no queda reflejado en el grafo. El ejemplo enfatiza la necesidad del uso de la temporalidad en la resolución de entidades.

\subsection{Graph Unlearning: eliminación de conocimiento en grafos}\label{subsec:graph-unlearning}

El \emph{graph unlearning} no es más que la aplicación del concepto de \emph{unlearning}, explicado en la sección~\ref{sec:graphrag-foundations}, a grafos. En ciertos contextos regulatorios surge la necesidad de eliminar de la base de conocimiento, de forma eficiente y segura, el efecto de datos concretos. La aplicación de las técnicas de \emph{graph unlearning} también es útil en otros escenarios donde el conocimiento del grafo debe eliminarse o actualizarse.

A continuación, se hace el análisis de tres resultados de investigación que incorporan el \emph{graph unlearning} en su concepción: GraphEraser \cite{graph_unlearning2022}, GIF \cite{gif2023} y ScaleGUN \cite{scalegun2025}.

GraphEraser aplica sobre el grafo un particionamiento balanceado. Para ello, se apoya en dos algoritmos: BLPA \cite{ugander2013blpa}, basado en propagación de etiquetas con restricciones de balance \cite{zhu2002labelpropagation}, y BEKM \cite{graph_unlearning2022}, una variante de \emph{k}-means \cite{kanungo2002kmeans} sobre el \emph{embedding} del grafo. Mediante el particionamiento se aíslan subgrafos independientes del grafo original. Una vez logrado el mismo, cuando se desea eliminar un dato del grafo (nodo o arista), basta con determinar en qué partición estaba y solo se reentrenar la misma.

Por su parte, GIF (\emph{Graph Influence Function}) cambia de estrategia con respecto a lo anteriormente descrito. Modela el impacto de eliminar un dato del grafo usando las dependencias estructurales hasta una distancia $k$. Esto se interpreta de la siguiente forma: la eliminación de un dato en el grafo repercute en los nodos y aristas vecinos que se encuentran a una determinada distancia de este. La relevancia de este hecho en los grafos de conocimiento puede interpretarse de la siguiente forma: eliminar una relación del KG afecta a las entidades, a las comunidades y a los resúmenes implicados en la misma.

ScaleGUN extiende las ideas que caracterizan a los sistemas anteriores para su aplicación en grafos con miles de millones de aristas. Usa una propagación local \emph{lazy} que pospone las actualizaciones sobre el grafo, entre ellas la eliminación de datos, hasta que sea necesario realizarlas. Solo actualiza los \emph{embeddings} afectados.

Los tres sistemas demuestran que la eliminación de conocimiento en grafos puede realizarse de manera eficiente, escalable y con garantías formales. No obstante, comparten una limitación semántica crítica en el contexto de grafos de conocimiento temporales: todos se basan exclusivamente en la operación de eliminación, sin considerar las causas que motivan la supresión de un dato. En muchos casos, en lugar de eliminar información, resulta más adecuado modificar la estructura del grafo. Por ejemplo, un hecho que ha dejado de ser verdadero, puede ser que el CEO de una determinada empresa ha cambiado. En este caso, esto representa una evolución en la plantilla de cuadros de la empresa, por tanto, el hecho existente en la base de conocimiento no debe eliminarse sino modificarse.

Por otra parte, un ejemplo de una situación en la que un hecho debe eliminarse del KG ocurre cuando este representa información sensible de la empresa que debe ser olvidada para preservar la privacidad de la misma. Este escenario se enmarca en el concepto de \emph{unlearning}; por tanto, en estos casos procede exclusivamente la eliminación del dato, junto con sus correspondientes efectos en cascada.

El aspecto esencial a destacar es que los sistemas estudiados reducen el concepto de \emph{unlearning} exclusivamente a la eliminación de datos. Este enfoque resulta insuficiente para un sistema que gestione de manera explícita la temporalidad en los grafos de conocimiento (KG).

%% ============================================================
\section{Benchmarks y Evaluación de Sistemas Temporales}\label{sec:benchmarks}

La evaluación de sistemas que integran grafos de conocimiento y razonamiento temporal requiere el uso de \emph{benchmarks} que aborden las distintas dimensiones del problema. En esta sección se hace un análisis crítico de diversas métricas representativas que han sido encontradas en el estudio del estado del arte.

El \emph{benchmark} CRAG \cite{crag2024}, ya introducido en la sección~\ref{sec:graphrag-foundations} del presente documento, es el más representativo para sistemas RAG con componentes temporales. En la mencionada investigación, dicho \emph{benchmark} fue aplicado sobre 4\,409 datos expresados en forma de pares (pregunta--respuesta) y un Mock KG de 2.6 millones de entidades. O sea, se aplica sobre ambos a la vez. CRAG, además, distingue de forma explícita la dificultad de determinar la respuesta a una pregunta según la velocidad de cambio del conocimiento. Define cuatro categorías de dinamismo para clasificar el tipo de preguntas:
\begin{itemize}
    \item tiempo real (\emph{real time}),
    \item cambio rápido (\emph{fast changing}),
    \item cambio lento (\emph{slow changing}),
    \item cambio estático (\emph{static change}).
\end{itemize}
Los resultados de la experimentación realizada muestran que las preguntas del tipo \emph{fast-changing} y \emph{real-time} son las más difíciles de responder en todos los sistemas evaluados.

El \emph{benchmark} introduce también \emph{mock Knowledge Graphs}. Estos permiten medir el efecto de la calidad del grafo sobre la calidad de la respuesta. Sin embargo, los \emph{mock KGs} se construyen una sola vez. El \emph{benchmark} no define un protocolo para actualizarlos cuando llegan documentos nuevos. En la práctica funcionan como \emph{snapshots} estáticos. CRAG mide la calidad de las respuestas temporales, pero no la del proceso de actualización del grafo.

Por su parte, el \emph{benchmark} TGB (\emph{Temporal Graph Benchmark}) \cite{tgb2023} evalúa qué tan bueno es un modelo al ser aplicado sobre grafos dinámicos. Su corpus reúne nueve \emph{datasets}, siete de los cuales fueron creados por ellos. TGB evalúa dos tareas sobre los modelos:
\begin{itemize}
    \item predicción de enlaces (\emph{link prediction}) \cite{libennowell2007link};
    \item afinidad de nodos (\emph{node affinity}) \cite{tgb2023}.
\end{itemize}
Ambas evaluaciones resultan útiles para sistemas basados en grafos temporales; sin embargo, no contemplan esquemas de evaluación del tipo pregunta--respuesta (Q/A), ni detección de conflictos temporales. Ambos son aspectos fundamentales en la concepción de un sistema GraphRAG temporal.

A modo de conclusión, puede señalarse que CRAG evalúa respuestas temporales sobre grafos estáticos, mientras que TGB aborda grafos dinámicos, aunque limitado a tareas de \emph{link prediction}. Sin embargo, ni estos enfoques ni otros identificados en la revisión del estado del arte contemplan de manera unificada: (a) la corrección temporal de las respuestas, (b) la detección y resolución de conflictos entre hechos contradictorios y (c) la calidad de la actualización incremental del grafo. La ausencia de un \emph{benchmark} que integre estas dimensiones, junto con la necesidad de su desarrollo, pone de manifiesto la importancia de incorporar al estado del arte un marco de evaluación que las aborde de forma conjunta.

%% ============================================================
\section{Conclusiones}\label{sec:soa-conclusions}

En el presente capítulo se realiza un análisis de los avances investigativos más recientes sobre el tratamiento de la temporalidad en sistemas de Generación Aumentada por Recuperación basados en grafos de conocimiento. Tras haber realizado dicho análisis se ha creado la Tabla~\ref{tab:state-of-the-art-comparison} en la cual se reflejan los aspectos técnicos que han sido identificados como relevantes para la concepción de un sistema de GraphRAG temporal. Para cada uno de los sistemas analizados, se expresa cuál es su propuesta con respecto a cada uno de dichos aspectos.

\begin{table}[htbp]
\centering
\begingroup
\scriptsize
\setlength{\tabcolsep}{3pt}
\renewcommand{\arraystretch}{1.15}
\begin{tabular}{@{}p{1.55cm}p{1.8cm}p{1.9cm}p{2.45cm}p{2.2cm}p{1.45cm}p{1.75cm}@{}}
\toprule
Sistema & Temporalidad en el grafo & Modelo de datos temporal & Estrategia ante conflictos & Resolución de entidades en actualización & Actualización incremental & Corrección vs Evolución \\
\midrule
GraphRAG \cite{edge2024local} & Claims con fechas y \texttt{status} & Intervalos en \emph{claims} & Coexistencia sin resolución & Fusión basada en coincidencia de nombre & Sí (\texttt{update}) & No distingue \\
T-GRAG \cite{tgrag2025} & Timestamps en hechos & Punto (tiempo válido) & Separación temporal (filtrado en consulta) & Extracción LLM sin alineación formal & Parcial & No distingue \\
TG-RAG \cite{tg_rag2025} & Aristas + jerarquía temporal & Punto (tiempo válido) & Separación temporal (aristas paralelas; filtrado) & No especificada & Sí (time reports selectivos) & No distingue \\
TREK \cite{trek2025} & Intervalos de validez & Intervalo (tiempo válido) & Detección + resolución por confianza & Embedding + LLM reranking & Sí (batch) & Parcial (excl./no-excl.) \\
DyG-RAG \cite{dyg_rag2025} & DEUs con timestamp & Punto (tiempo válido) & No contempla & Coreference resolution & Inserción incremental & No distingue \\
E2RAG \cite{e2rag2025} & Eventos con orden narrativo & Implícito (mención-evento) & No deduplica (preserva versiones) & Evitada (sin deduplicación) & Parcial & No distingue \\
ZEP/Graphiti \cite{zep2025} & Intervalos bitemporales & Bitemporal (válido + transacción) & Detección por similitud semántica + LLM & Búsqueda híbrida + LLM & Sí (streaming) & Parcial (invalidación) \\
STAR-RAG \cite{star_rag2025} & Reglas con timestamps & Punto (tiempo válido) & No contempla & No especificada & No (batch) & No distingue \\
\bottomrule
\end{tabular}
\endgroup
\caption{Comparación de sistemas según dimensiones técnicas del análisis.}
\label{tab:state-of-the-art-comparison}
\end{table}

A partir de la información reflejada en la tabla se identifican seis limitaciones que comparten los sistemas analizados, las cuales se mencionan y describen a continuación:
\begin{enumerate}
    \item \textbf{GraphRAG: temporalidad presente, pero no integrada en el \emph{pipeline}.} En GraphRAG existen las marcas temporales en los \emph{claims} así como la actualización incremental del sistema, pero ninguna de las dos se integra con los procesos centrales del \emph{pipeline}: no filtran la recuperación, no intervienen en la detección de comunidades y no afectan la descripción de las comunidades.

    \item \textbf{Separación temporal en lugar de resolución de conflictos.} Los sistemas más maduros filtran los hechos por fecha en la consulta y así evitan que los que son contradictorios coincidan en la respuesta. La estrategia falla cuando la contradicción se da dentro del mismo período de tiempo, cuando no hay una fecha clara que separe lo nuevo de lo anterior, o cuando hace falta distinguir un error del pasado de una evolución en los hechos.

    \item \textbf{Cardinalidad de relaciones poco desarrollada.} La única etiqueta existente (exclusiva / no exclusiva, en TREK) es simétrica y no captura que la restricción pueda aplicarse solo en uno de los extremos de la relación, lo que genera falsos positivos y pasa por alto conflictos reales.

    \item \textbf{Bitemporalidad fuera del \emph{pipeline} de GraphRAG.} El modelo bitemporal (distinguir cuándo un hecho fue verdadero en el mundo y cuándo el sistema lo registró) está consolidado en teoría y validado en sistemas dentro de otros contextos \cite{zep2025}, pero no se ha integrado con la sumarización jerárquica y la búsqueda por comunidades de GraphRAG.

    \item \textbf{Resolución de entidades sin considerar la temporalidad.} Los sistemas analizados ante la presencia de entidades duplicadas resuelven dicho conflicto mediante la similitud entre nombres o entre \emph{embeddings}. No consideran que dos menciones, en períodos de tiempo distintos, pueden describir estados distintos de una misma entidad.

    \item \textbf{Corrección y evolución sin distinción formal.} Cuando un hecho nuevo se inserta en el grafo de conocimiento y se relaciona con uno ya existente, los sistemas actuales aplican un mismo procedimiento de actualización (por lo general, sobrescribir el hecho anterior o añadir el nuevo en paralelo) sin diferenciar tres situaciones que son semánticamente distintas: un hecho que dejó de ser verdadero (evolución), un hecho que nunca lo fue y se reconoce como error (corrección) y un hecho nuevo que confirma uno ya presente (corroboración).
\end{enumerate}

Estas limitaciones están relacionadas entre sí. La ausencia de un modelo bitemporal (limitación 4) impide distinguir entre corrección y evolución (limitación 6). La falta de sensibilidad temporal en la resolución de entidades (limitación 5) impide una detección de conflictos precisa (limitación 2). La escasa expresividad de la clasificación de la cardinalidad de las relaciones (limitación 3) restringe la dirección desde la cual buscar esos conflictos. La limitación 1, especialmente, está relacionada con todas las demás, ya que, aunque existen sistemas que han resuelto algunas de ellas, dicha solución no se ha integrado al \emph{pipeline} de GraphRAG.

Estas seis limitaciones sustentan la pertinencia y viabilidad de la presente investigación en función de los objetivos propuestos. Asimismo, constituyen la base para las decisiones de diseño de la solución presentada en el capítulo siguiente. Dicha propuesta se desarrolla dentro del \emph{pipeline} de GraphRAG ---manteniendo la sumarización jerárquica y la búsqueda por comunidades--- e integra cuatro componentes que responden directamente a resolver las limitaciones identificadas:
\begin{itemize}
    \item[(a)] un modelo bitemporal a nivel de arista (limitaciones 4 y 6, y, en consecuencia, la 1);
    \item[(b)] un mecanismo de resolución de conflictos (limitación 2);
    \item[(c)] un esquema de clasificación de cardinalidad asimétrica (limitación 3);
    \item[(d)] un mecanismo de resolución temporal de entidades (limitación 5).
\end{itemize}
